\documentclass[11pt]{article}
\usepackage[a4paper,margin=21mm]{geometry}
\usepackage[no-math]{fontspec}
\setmainfont{lmroman10-regular.otf}[BoldFont=lmroman10-bold.otf,ItalicFont=lmroman10-italic.otf,BoldItalicFont=lmroman10-bolditalic.otf]
\usepackage{amsmath,amssymb,amsthm,amscd,graphicx,color}
\usepackage{xurl}
\usepackage[unicode,hidelinks]{hyperref}
\setlength\emergencystretch{3em}
\numberwithin{equation}{section}
\newtheorem{theorem}{Theorem}[section]
\newtheorem{lemma}[theorem]{Lemma}
\newtheorem{proposition}[theorem]{Proposition}
\newtheorem{remark}[theorem]{Remark}
\allowdisplaybreaks
\begin{document}
\begin{center}\Large General energy decay for a coupled viscoelastic wave system with a single memory term\end{center}
\noindent\textbf{Stable ID:} 3413\par
\noindent\textbf{Author:} Houda Fadel; Djamila Benterki; Salim A. Messaoudi\par
\noindent\textbf{Work:} Existence and stability of solutions for a viscoelastic coupled system of two wave equations\par
\noindent\textbf{Source:} \url{https://ejde.math.txstate.edu/Volumes/2026/16/abstr.html}\par
\noindent\textbf{Version:} Electronic Journal of Differential Equations, Vol. 2026 (2026), No. 16, pp. 1--23; published February 19, 2026; official native/PDF version of record acquired 2026-10-04\par
\noindent\textbf{Rights:} CC BY 4.0. \url{https://creativecommons.org/licenses/by/4.0/}. Independent typesetting; source mathematics preserved.\par
\noindent\textbf{Difficulty:} H2. Only the first wave is directly damped by memory, so positivity of the total energy alone does not control the undamped second wave or yield a general kernel-dependent decay rate. The authors assemble several cross and memory perturbations into a Lyapunov functional, balance coupling constants under alpha<sqrt(l)/Cp, then use both first- and second-order energy to derive a nonlinear convexity bound for the memory tail. Converting that bound into the inverse-G2 rate for a general nonlinear relaxation function is a nonroutine residual stability argument beyond standard wave energy estimates.\par
\noindent\textbf{Editorial scope:} One original published theorem unit: Theorem4.10 (native label thm13); all original A1-A3, regularity, constants, time quantifiers, linear/nonlinear alternatives and inverse-G2 conclusions unchanged. Full weak-solution construction and all first/second-order energy, cross-memory and Lyapunov estimates included. Numerical examples unselected; no branch or parameter variant counted separately.\par
\medskip\hrule\medskip
% Source-native exact authored range 73--223
 \section{Introduction}
 
 Wave equations with memory effects have been a central topic in the analysis of partial
 differential equations, owing to their applications in physics, engineering, and material
 science. The focus has been on the existence, stability, and energy decay of solutions, with significant progress made for both single and coupled wave equations.
 To motivate our work, data collected and results are reviewed to provide a thorough study of
 both the single and coupled viscoelastic wave equation. We start with the pioneer works of Dafermos \cite{Dafermos1, Dafermos2}, where the author showed that solutions decay to zero
 for smooth, monotone relaxation functions $g$, but the decay rate remained unspecified. Cavalcanti et al.\ \cite{CavalcantiDS} studied
\[
|u_t|^\rho u_{tt} - \Delta u - \Delta u_{tt} + \int_0^t g(t - s) \Delta u(s) \, ds - \gamma \Delta u_t = 0,\quad \text{in } \Omega \times \mathbb{R}^{*}_{+}
\]
proving global existence for $\gamma \geq 0$ and exponential energy decay for $\gamma > 0$. Messaoudi and Tatar \cite{MTatar1, MTatar2} extended these results to include source terms, showing exponential and polynomial decay depending on the rate of decay of g(t).
Cavalcanti et al.\ \cite{CavalcantiFerreira} considered a viscoelastic wave equation with a
local frictional damping, where the kernel $g$ satisfies, for two positive constants $\xi_1$
and $\xi_2$,
\begin{align*}
-\xi_1 g(t) \leq g'(t) \leq -\xi_2 g(t), \quad \forall t \in \mathbb{R}_{+}.
\end{align*}
Under the above assumptions and subject to certain restrictions on the control zone, they established an exponential decay result. Berrimi and Messaoudi \cite{Berrimi2004} later
improved this result by proving that the viscoelastic dissipation alone is sufficient to stabilize the system, and subsequently extended it to the case where a source term competes with the dissipation \cite{Berrimi2006}.
Further contributions including the analysis of nonlinear damping with polynomially decaying kernels \cite{Cavalcanti2003} and global existence with uniform decay for source-type problems \cite{Li2011}  appeared later.
 Without restricting $g(t)$ to exponential or polynomial decay, Messaoudi \cite{SMessaoudi1, SMessaoudi2} generalized all the existing decay results by considering the following equation
\[
u_{tt} - \Delta u + \int_0^t g(t - \tau) \Delta u(\tau) \, d\tau = b |u|^{\gamma} u,\quad b \geq 0,
\]
where $g'(t) \leq -\xi(t) g(t)$, for a non-increasing differentiable $\xi(t)$.
These latter results inspired other authors and more results have been established, see \cite{SaidH2011, Hao, PJS, WSH, MessaoudiMustafa}.

Further improvement was establish by Messaoudi and Al-Khulaifi \cite{MessaoudiWaled}, who addressed the problem in \cite{CavalcantiDS}, for $\gamma =0$ and for a kernel satisfying, for $1\leq p\leq \frac{3}{2}$, the condition
\[
g'(t) \leq -\xi(t) g^p(t), \quad \forall t \in \mathbb{R}_{+}.
\]
They proved that the energy of solutions converges to zero at the same rate as $g(t)$, provided $g(t)$ decays faster than $s^{-2}$ at infinity. Later, Mustafa \cite{MM181} extended his earlier joint work with Messaoudi \cite{MM}, which dealt with kernels satisfying $g'(t) \leq -G(g(t))$, to the more general condition $$g'(t) \leq -\xi(t) G(g(t)), \quad \forall t \in \mathbb{R}_{+},$$ where $G$ is strictly increasing and convex function. This extension provided explicit energy decay estimates and encompassed a wider range of decay behaviors.

Coupled wave systems also have consistently drawn attention due to their fundamental role in understanding stability and long-term dynamics. In a recent contribution, Guesmia, Messaoudi, and Zahri \cite{GASM} analyzed an abstract coupled system involving a viscoelastic component interacting with a purely elastic one. Precisely, they looked into the problem
\begin{gather*}
u_{tt}(t) + Au(t) - \int_0^{t} g(s) A^{\theta} u(t - s) \, ds + \alpha u + \beta B v(t) = f(u), \quad t > 0, \\
v_{tt}(t) + A v(t) + \beta B u(t) = 0, \quad t > 0,\\
u(0) = u_0, \quad u_t(0) = u_1, \quad v(0) = v_0, \quad v_t(0) = v_1,
\end{gather*}
 where $A$ is a self-adjoint positive operator on a Hilbert space $H$, $\theta \in [0,1]$,
 $\alpha > 0$, and $g$ is a positive nonincreasing relaxation function.
 Their analysis established general decay rates for the system energy, showing that the
 asymptotic behavior depends strongly on the relaxation function and the coupling parameter.
 Numerical experiments were also provided, highlighting the impact of memory effects on
 the qualitative behavior of solutions. Earlier, Feng and Li \cite{FengBaoweiHaiyan} considered
 a viscoelastic Lam\'e system with strong damping terms. For two distinct relaxation kernels
 $g_1$ and $g_2$, they assumed conditions of the form $g_i'(t) \leq -\xi_i(t) H_i(g_i(t))$
 ($i=1,2$), which led to explicit and general decay results. Their work improved earlier results
 of Beniani, Taouaf, and Benaissa \cite{Beniani}, who had previously established a general
 decay property for similar coupled viscoelastic structures.
Further investigations were carried out by Jin et al. \cite{JKJT}, who studied a coupled system of two abstract evolution equations with finite memory, given by
\begin{gather*}
 u_{tt}(t) + Au(t) - \int_0^{t} g(s) A u(t - s) \, ds + \alpha u + \beta B v(t)
 = f(u), \quad t > 0, \\
v_{tt}(t) + A v(t) + \beta B u(t) = 0, \quad t > 0,
 \end{gather*}
 where $A$ and $B$ are operators, $\alpha \geq 0$, $\beta > 0$, and $f$ is a nonlinear source.
 Under appropriate assumptions on $g$ and structural conditions on $A$, $B$, and $f$,
 they proved polynomial stability with decay rates of order $t^{-1}$. Several extensions and
related contributions in this direction may be found in \cite{ ORHC, AkilM}.
In addition to viscoelastic systems, purely elastic coupled wave models have also been studied.
For instance, a weakly coupled system
\begin{gather*}
u_{tt} - \Delta u + u_t + \alpha v = 0, \quad \text{in } \Omega \times (0, \infty), \\
v_{tt} - \Delta v + \alpha u = 0, \quad \text{in } \Omega \times (0, \infty),
\end{gather*}
was analyzed in \cite{LRF}, where $\Omega \subset \mathbb{R}^n$ is a bounded domain and
$\alpha$ a positive constant. The authors showed that the energy decays polynomially,
with the rate being optimal, and provided numerical validation in the one-dimensional case.
Similar results on stability and decay rates had earlier been obtained for systems with
boundary damping, notably by Komornik et al. \cite{KVBR} and Aassila \cite{Assila1, Assila2},
where exponential stability was achieved in the case of linear damping and polynomial stability
in the case of nonlinear damping mechanisms.

Motivated by the findings in \cite{GASM}, the present work investigates the coupled viscoelastic wave system
\begin{equation} \label{e1.1}
\begin{gathered}
u_{tt} - \Delta u + \int_0^t g(t - s) \Delta u \, ds + \alpha v = 0, \quad
 \text{in } \Omega \times (0, T), \\
v_{tt} - \Delta v + \alpha u = 0, \quad \text{in } \Omega \times (0, T), \\
u = v = 0, \quad \text{on } \partial \Omega \times (0, T), \\
u(0) = u_0, \, u_t(0) = u_1, \; v(0) = v_0, \, v_t(0) = v_1, \quad \text{in } \Omega,
\end{gathered}
\end{equation}
where \( \Omega \) is a bounded domain in \( \mathbb{R}^N, \) \( \alpha > 0 \), and the initial
data belong to suitable spaces. The relaxation function \( g \) satisfies
\( g'(t) \leq -\xi(t) G(g(t)) \), where $G$ is an increasing and convex function near the origin and
$\xi$ is a nonincreasing function. We establish, in details, an existence and uniqueness result,
then we prove a general and explicit result for energy decay and provide some numerical
illustration to validate our theoretical findings.

The structure of this paper is as follows: the next section presents some assumptions,
remarks and functionals needed in our results. In Section 3, we establish the existence of
solutions by combining energy estimates with the Faedo-Galerkin method.
Section 4 presents several technical lemmas and contains the statement and proof of the
decay results.
In section 5, two numerical tests are conducted aiming to verify our decay result.

 \section{Preliminaries}

 This section presents the foundation for our main result. We set
 \begin{gather*}
 (h \star \varphi)(t) = \int_0^t h(t-s) \varphi(s)ds,\
 (h \diamond \varphi)(t) = \int_0^t h(t-s)\| \varphi(t)- \varphi(s)\| ds, \\
 (h \circ \varphi)(t) = \int_0^t h(t-s)\int_\Omega\| \varphi(t)- \varphi(s)\|^2 dx \, ds.
 \end{gather*}
 For the relaxation function $g$, we assume:
 \begin{itemize}
\item[(A1)] $g: \mathbb{R}^+ \to \mathbb{R}^+ $ is a $ C^1$ function satisfying\\
\begin{equation*}
 g(0) > 0, \quad 1-\int_0^{\infty} g(s) \, ds = l > 0. \label{e2.1}
\end{equation*}

\item[(A2)] There exists a $C^1$ function $G :\mathbb{R}^+ \to \mathbb{R}^+ $
which is linear or it is strictly increasing and strictly convex $C^2$ function on
$(0, r]$, $r \leq g(0)$, with $G(0) = G'(0) = 0$, such that
\begin{equation} \label{e2.2}
g'(t) \leq -\xi(t) G(g(t)),\quad \forall t \geq 0,
\end{equation}
where $\xi$ is a positive non-increasing differentiable function.

\item[(A3)] The constant $\alpha $ is such that
\begin{equation} \label{e2.3}
0 < \alpha < \frac{\sqrt{l} }{C_p},
\end{equation}
where $C_p$ is the Poincar\'e constant.
\end{itemize}

\begin{lemma}[\cite{Munoz}] \label{lem1}
For any function $h \in C^1 (\mathbb{R})$ and any $k \in H^1 (0,T)$, we have
$$
(h \star k)(t) k_t
 = -\frac{1}{2} h(t)  \| k\|^2 + \frac{1}{2} (h' \circ k)(t)
 -\frac{1}{2} \frac{d}{dt} \Big\{( h \circ k)(t)
 - \Big( \int_0^t h(\tau) d\tau \Big) \|k\|^2 \Big\}.
$$
\end{lemma}

\begin{lemma}[\cite{SMessaoudi1}] \label{lem2}
 For $u \in H^1_0(\Omega)$, we have
\begin{equation}
\int_\Omega \Big( \int_0^t g(t - s)(u(t) - u(s)) \, ds \Big)^2 dx
\leq (1 - l) C_p^2 (g \circ \nabla u)(t),\label{e2.4}
\end{equation}
where $l$ is given in {\rm (A1)}.
\end{lemma}

\section{Existence of weak solutions}
Now, we examine the well-posedness of problem \eqref{e1.1}. To establish this, we will combine some energy estimates and the Faedo-Galerkin method.

% Source-native exact authored range 224--633
\begin{theorem} \label{thm3}
 Assume that {\rm (A1)--(A3)} hold.
Then, for each $ (u_0,v_0) \in \left( H^2(\Omega) \cap H^1_0(\Omega)\right)^2$ and
$(u_1,v_1) \in \left( H^1_0(\Omega)\right)^2$, problem \eqref{e1.1} has a unique strong solution
$(u, v)$ satisfying, for every $T > 0$,
\begin{gather*}
u, v \in L^{\infty}\left( [0, T), H^2(\Omega) \cap H^1_0(\Omega)\right),\\
u_{t}, v_{t} \in L^2\left( (0, T), H^1_0(\Omega)\right)\cap H^1_0 \left( (0, T)\times \Omega\right),\\
u_{tt}, v_{tt} \in L^{\infty}\left( (0, T), L^2(\Omega)\right).
\end{gather*}
 \end{theorem}

\begin{proof} \quad

\noindent\textbf{Step 1.} Faedo-Galerkin approximation.
We take
\(\{ \phi_j \}_{j=1}^\infty\)
to be the eigenfunctions of the Laplacian operator subject to Dirichlet boundary conditions, satisfying $-\Delta \phi_i = \lambda_i \phi_i $. Then
$\{\phi_j\}_{j=1}^\infty$ forms an orthogonal basis of $ H_0^1(\Omega)$,
\( H^2(\Omega) \cap H_0^1(\Omega) \), and \( L^2(\Omega) \).
Let $V_{N}$ be spanned by the $N$ first functions $ \phi_1, \phi_2, \dots, \phi_N $.
We seek  an approximate solution $(u,v)$ of the form
\[
u^N(x,t) = \sum_{j=1}^N a_j(t) \phi_j(x), \quad
v^N(x,t) = \sum_{j=1}^N b_j(t) \phi_j(x), \quad t\in (0,T),
\]
which satisfies
\begin{gather*}
\int_\Omega u_{tt}^N \phi_j dx + \int_\Omega \nabla u^N .\nabla \phi_j dx - \int_\Omega \int_0^t g(t-s) \nabla u^N(s) .\nabla \phi_j\, ds\,dx+ \alpha \int_\Omega v^N \phi_j \, dx = 0,\\
\int_\Omega v_{tt}^N \phi_j dx + \int_\Omega \nabla v^N .\nabla \phi_j dx+ \alpha \int_\Omega u^N \phi_j \, dx = 0,\quad 1\leq j\leq N
\label{e3.1}
\end{gather*}
and
\begin{equation} \label{e3.2}
\begin{gathered}
 u^{N}(0) = u^{N}_0= \sum_{i=1}^{N} \langle u_0, \phi_i \rangle \phi_i
 \to u_0, \\
  v^{N}(0) = v^{N}_0 = \sum_{i=1}^{N} \langle v_0, \phi_i \rangle \phi_i
 \to v_0 \quad \text{in } H^2(\Omega) \cap H_0^1(\Omega), \\
 u_{t}^N(0) = u_1^N = \sum_{i=1}^{N} \langle u_1, \phi_i \rangle \phi_i \to u_{1}, \\
 v_{t}^N(0) = v_1^N = \sum_{i=1}^{N} \langle v_1, \phi_i \rangle \phi_i \to v_{1} \quad \text{in } H^1_0(\Omega).
 \end{gathered}
\end{equation}
Standard results of ordinary differential equations guarantee the existence of
a unique local solution
$ (u^N, v^N )$ of \eqref{e3.1}-\eqref{e3.2} defined on
$ [0, t_{N}), 0 < t_{N} \leq T $.
Our next step is to obtain a priori estimates that enable us to extend the
solution to the entire interval.
\smallskip

\noindent\textbf{Step 2.} A priori estimate 1.
By multiplying the first equation in \eqref{e3.1} by \(a'_j(t)\) and summing the results over \(j\)
from 1 to \(N\), we obtain, for all \(0 < t \leq t_N\),
\begin{equation}
\begin{aligned}
&\frac{d}{dt} \Big( \frac{1}{2} \| u^N_{t}(t) \|^2 + \frac{1}{2} \| \nabla u^N(t) \|^2 \Big)
+ \alpha \int_{\Omega} v^N(t) \, u^N_{t}(t) \, dx
&- \int_0^t g(t-s) \int_{\Omega} \nabla u^N(s) \nabla u^N_{t}(t) \, dx ds =0.
\end{aligned}\label{e3.3}
\end{equation}
Using Lemma \ref{lem1} for the last term of the left-hand side of \eqref{e3.3}
we find that
\begin{equation}
\begin{aligned}
&\frac{d}{2dt} \Big( \| u^N_{t}(t) \|^2
+ \Big[ 1-\int_0^t g(s) ds \Big] \| \nabla u^N(t) \|^2 + (g\circ \nabla u^N)(t) \Big)
+\alpha \int_{\Omega} v^N(t) \, u^N_{t}(t) dx \\
& +\frac{1}{2} g(t)\parallel \nabla u^N(t) \parallel^2 - \frac{1}{2}( g'\circ\nabla u^N) (t) =0.
\end{aligned} \label{e3.4}
 \end{equation}
Multiplying the second equation in \eqref{e3.1} by $b'_j(t)$ and summing with respect to $j$, we obtain
 \begin{equation}
 \frac{d}{2 dt} \Big( \| v^N_{t}(t) \|^2 + \| \nabla v^N (t) \|^2 \Big)
 + \alpha \int_{\Omega} u^N(t) \, v^N_{t}(t) \, dx =0, \quad 0 < t \leq t_N.
 \label{e3.5}
\end{equation}
Then, combining \eqref{e3.4} and \eqref{e3.5} yields
\begin{equation}
\begin{aligned}
 &\frac{d}{2 dt} \Big( \| u^N_{t}(t) \|^2 + \Big[ 1-\int_0^t g(s)ds\Big]
 \| \nabla u^N(t) \|^2 +(g\circ\nabla u^N)(t)+\| v^N_{t}(t) \|^2 + \| \nabla v^N(t) \|^2 \\
 &+ 2 \alpha \int_{\Omega} v^N(t) u^N(t) dx\Big) \\
&=-\frac{1}{2} g(t)\parallel \nabla u^N(t) \parallel^2 + \frac{1}{2}( g'\circ\nabla u^N )(t).
\end{aligned}\label{e3.6}
\end{equation}
By assumptions (A1)--(A3) we obtain
 \begin{equation}
\begin{aligned}
&\frac{d}{2 dt} \Big( \| u^N_{t} (t) \|^2 + \Big[ 1-\int_0^t g(s)ds\Big]
 \| \nabla u^N(t) \|^2 +(g\circ\nabla u^N)(t)+\| v^N_{t}(t) \|^2
 + \| \nabla v^N (t)\|^2 \\
&+ 2 \alpha \int_{\Omega} v^N(t) u^N(t) dx\Big)\leq 0.
\end{aligned} \label{e3.7}
\end{equation}
Integrating \eqref{e3.7} over $(0, t)$ yields
\begin{equation}
 E_N(t)\leq E_N(0), \label{e3.8}
\end{equation}
where
 \begin{gather*}
 \begin{aligned}
E_N(t)&=\| u^N_{t} (t) \|^2 + \Big[ 1-\int_0^t g(s)ds\Big] \| \nabla u^N(t) \|^2 +(g\circ\nabla u^N)(t)
 +\| v^N_{t}(t) \|^2 + \| \nabla v^N(t) \|^2\\
&\quad + 2 \alpha \int_{\Omega} v^N(t) u^N(t) dx,
\end{aligned}\\
E_N(0)= \| u^N_{1} \|^2 + \| v^N_{1} \|^2 +\| \nabla u^N_0 \|^2
 +\| \nabla v^N_0 \|^2+ 2 \alpha \int_{\Omega} v^N_0 u^N_0 dx.
\end{gather*}
Using Young's and Poincar\'e's inequalities, for $\varepsilon_{1}> 0$, we have
 \begin{equation}
 2\int_{\Omega} v^N(t) u^N(t) dx\geq - \varepsilon_{1} C_p\| \nabla u^N (t) \|^2 - \frac{C_p}{\varepsilon_{1}} \| \nabla v^N (t) \|^2
 \label{e3.9}
 \end{equation}
and
 \begin{equation}
 2 \alpha \int_{\Omega} v^N_0 u^N_0 dx\leq \alpha C_p\| \nabla u^N_0 \|^2 + \alpha C_p \| \nabla v^N_0 \|^2.
 \label{e3.10}
 \end{equation}
Thus
\begin{equation}
 E_N(0)\leq \| u^N_{1} \|^2 + \| v^N_{1} \|^2 +(1+\alpha C_p)\| \nabla u^N_0 \|^2 +(1+\alpha C_p)\| \nabla v^N_0 \|^2.\label{e3.11}
 \end{equation}
Since the sequences $ (u_0^N)_{N \in \mathbb{N}}$,
$ (u_1^N)_{N \in \mathbb{N}}, (v_0^N)_{N \in \mathbb{N}}, (v_1^N)_{N \in \mathbb{N}} $ converge, there exists a positive constant $L_{1}$, independent of $N$, such that
\begin{equation}
E_N(0)\leq L_{1}.\label{e3.12}
\end{equation}
On the other hand,
\begin{equation}
\begin{aligned}
E_N(t)
&\geq \| u^N_{t}(t) \|^2 + \Big[ 1-\int_0^t g(s)ds -\alpha \varepsilon_{1} C_p \Big]
 \| \nabla u^N(t) \|^2 +(g\circ\nabla u^N)(t)+\| v^N_{t} (t) \|^2 \\
 &\quad +\Big( 1-\frac{\alpha C_p}{\varepsilon_{1}}\Big) \| \nabla v^N(t) \|^2.
\end{aligned}\label{e3.13}
\end{equation}
Let us select $\varepsilon_{1}$ such that it satisfies
$\alpha C_p \leq \varepsilon_{1}\leq \frac{l}{\alpha C_p}$.
Then, there exists $C>0$, such that
\begin{equation}
E_N(t)\geq C\left( \| u^N_{t}(t) \|^2 + \| \nabla u^N (t)\|^2
+(g\circ\nabla u^N)(t)+\| v^N_{t}(t) \|^2 + \| \nabla v^N (t)\|^2\right).\label{e3.14}
\end{equation}
From \eqref{e3.8} and \eqref{e3.12}, we deduce
\begin{equation}
\| u^N_{t}(t) \|^2 +\| v^N_{t}(t) \|^2 + \| \nabla u^N(t) \|^2
+ \| \nabla v^N (t)\|^2 +(g\circ\nabla u^N)(t)\leq \frac{L_{1}}{C},\quad
\forall t\leq t_{N}\leq T,
\label{e3.15}
\end{equation}
where $\frac{L_{1}}{C}$ is a constant independent of $t$ and $N$.
\smallskip

\noindent\textbf{A priori estimate 2.}
In \eqref{e3.1}, we substitute $\phi_j $ by $ -\Delta \phi_j$ then
multiply the equation \eqref{e3.1}$_1$ by $a'_j(t)$ and the second equation
\eqref{e3.1}$_2$ by $b'_j(t)$, sum the resulting equations, and use
(A1)--(A3), we arrive at
\begin{align*}
&\frac{d}{2dt}\Big[ \| \nabla u^N_{t}(t) \|^2
+\Big( 1 - \int_0^t g(s)\,ds \Big) \| \Delta u^N(t) \|^2
+ (g \circ \Delta u^N)(t) + \| \nabla v^N_{t}(t) \|^2
+ \| \Delta v^N(t) \|^2 \\
& +2 \alpha \int_{\Omega} \nabla v^N(t) \nabla u^N(t) dx\Big]
-\frac{1}{2} \left( g' \circ \Delta u^N \right)(t)
+\frac{1}{2} g(t) \| \Delta u^N(t) \|^2\leq 0.
\end{align*}

Applying Young's and Poincar\'e's inequalities, we have for $\varepsilon_{3}> 0$,
 \begin{equation*}
 2\int_{\Omega} \nabla v^N(t) \nabla u^N(t) dx
 \geq - \varepsilon_{3} C_p\| \Delta u^N(t) \|^2 - \frac{C_p}{\varepsilon_{3}}\| \Delta v^N(t) \|^2,
 \end{equation*}
and
\[
 2 \int_{\Omega} \nabla v^N_0 \nabla u^N_0 dx
 \leq C_p\| \Delta u^N_0 \|^2 +C_p \| \Delta v^N_0 \|^2.
\]
Then
\begin{equation}
\begin{aligned}
& \| \nabla u^N_{1} \|^2+ \| \nabla v^N_{1} \|^2+\| \Delta u^N_0 \|^2
 + \| \Delta v^N_0 \|^2 +2 \alpha \int_{\Omega} \nabla v^N_0 \nabla u^N_0 dx \\
&\leq \| \nabla u^N_{1} \|^2+ \| \nabla v^N_{1} \|^2
 +(1+\alpha C_p) \| \Delta u^N_0 \|^2
 +(1+\alpha C_p)\| \Delta v^N_0 \|^2.
\end{aligned} \label{e3.16}
\end{equation}
Since the sequences $ (u_0^N)_{N \in \mathbb{N}}$,
$ (u_1^N)_{N \in \mathbb{N}}, (v_0^N)_{N \in \mathbb{N}}, (v_1^N)_{N \in \mathbb{N}} $
converge, there exists a positive constant $L_2$, independent of $N$, such that
\begin{equation}
 \| \nabla u^N_{1} \|^2+ \| \nabla v^N_{1} \|^2
 +\| \Delta u^N_0 \|^2 + \| \Delta v^N_0 \|^2
 +2 \alpha \int_{\Omega} \nabla v^N_0 \nabla u^N_0 dx\leq L_2. \label{e3.17}
\end{equation}
Furthermore,
\begin{align*}
&\| \nabla u^N_{t}(t) \|^2+\Big( 1 - \int_0^t g(s)\,ds \Big) \| \Delta u^N(t) \|^2
 + (g \circ \Delta u^N)(t) + \| \nabla v^N_{t}(t) \|^2+ \| \Delta v^N (t)\|^2 \\
&+2 \alpha \int_{\Omega} \nabla v^N (t)\nabla u^N(t) dx \\
&\geq \| \nabla u^N_{t}(t) \|^2+\Big( 1 - \int_0^t g(s)\,ds
 - \varepsilon_{3} C_p \Big) \| \Delta u^N(t) \|^2 \\
&\quad + (g \circ \Delta u^N)(t) + \| \nabla v^N_{t}(t) \|^2
 +\Big( 1-\frac{\alpha C_p}{\varepsilon_{3}}\Big) \| \Delta v^N(t) \|^2.
\end{align*}

We choose $\varepsilon_{3}$ such that it satisfies
$\alpha C_p \leq \varepsilon_{3}\leq \frac{l}{\alpha C_p}$.
Then there exists $C>0$, such that
\begin{equation}
\begin{aligned}
&\| \nabla u^N_{t}(t) \|^2+\Big( 1 - \int_0^t g(s)\,ds \Big) \| \Delta u^N(t) \|^2
+ (g \circ \Delta u^N)(t) + \| \nabla v^N_{t}(t) \|^2+ \| \Delta v^N(t) \|^2\\
&+2 \alpha \int_{\Omega} \nabla v^N (t)\nabla u^N (t) dx \\
&\geq C \Big( \| \nabla u^N_{t}(t) \|^2+ \| \Delta u^N (t)\|^2
+ (g \circ \Delta u^N)(t) + \| \nabla v^N_{t}(t) \|^2
+ \| \Delta v^N (t)\|^2 \Big).
\end{aligned} \label{e3.18}
\end{equation}
From \eqref{e3.16}, \eqref{e3.17} and \eqref{e3.18}, it follows that
\begin{equation}
\| \nabla u^N_{t}(t) \|^2+ \| \Delta u^N(t) \|^2 + (g \circ \Delta u^N)(t)
+ \| \nabla v^N_{t}(t) \|^2+ \| \Delta v^N(t) \|^2
\leq \frac{L_2}{C}, \; \forall\, t\leq t_{N} \leq T,\label{e3.19}
 \end{equation}
where $\frac{L_2}{C}$ is a constant independent of $t$ and $N$.
\smallskip

\noindent\textbf{A priori estimate 3.}
 By multiplying the first equation in \eqref{e3.1}$_1$ by
 \(a''_j(t)\) and summing the results over \(j\) from 1 to \(k\), we obtain
\begin{align*}
&\int_\Omega |u_{tt}^N(t)|^2 dx \\
& = \int_\Omega \Delta u^N(t) u_{tt}^N(t) \, dx
 - \int_\Omega \int_0^t g(t-s) \Delta u^N(s) u_{tt}^N(t) \, ds\,dx
 - \alpha \int_\Omega v^N(t) u_{tt}^N(t) \, dx\\
& = \int_\Omega \Delta u^N(t) u_{tt}^N(t) \, dx
 + \int_\Omega u_{tt}^N(t) \int_0^t g(t-s) \Big(\Delta u^N(t)- \Delta u^N(s)\Big)\, ds\,dx\\
&\quad -\int_0^t g(s) ds \int_\Omega u_{tt}^N(t) \Delta u^N(t) dx
 - \alpha \int_\Omega v^N(t) u_{tt}^N(t) \, dx\\
&\leq \varepsilon \int_\Omega |u_{tt}^N(t)|^2 dx
 +\frac{1}{4\varepsilon}\int_\Omega |\Delta u^N(t)|^2 \, dx
 + \varepsilon \int_\Omega |u_{tt}^N(t)|^2 dx+\frac{1-l}{4\varepsilon}
 (g \circ \Delta u^N)(t)\\
&\quad +\varepsilon(1-l) \int_\Omega |u_{tt}^N(t)|^2 dx
 + \frac{1-l}{4\varepsilon} \int_\Omega |\Delta u^N(t)|^2dx
 + \varepsilon \int_\Omega |u_{tt}^N(t)|^2 dx
 +\frac{\alpha^2}{4\varepsilon}\int_\Omega |v^N(t)|^2 dx \\
&\leq \varepsilon (4-l)\int_\Omega |u_{tt}^N(t)|^2 dx
 +\frac{2-l}{4\varepsilon} \int_\Omega |\Delta u^N(t)|^2dx+
 \frac{1-l}{4\varepsilon} (g \circ \Delta u^N)(t)
 +\frac{\alpha^2}{4\varepsilon}\int_\Omega |v^N(t)|^2 dx.
\end{align*}
Hence, for some $c >0$, we have
\begin{align}
\left(1-\varepsilon(4-l)\right) \int_\Omega |u_{tt}^N(t)|^2 dx
\leq \frac{c}{4\varepsilon}(L_1+L_2).
\end{align}
We take $\varepsilon >0$ small enough, so
\begin{align}
\|u_{tt}^N(t)\|^2 \leq \tilde{c}(L_1+L_2)=L_3. \label{e3.20}
\end{align}
Similarly, we obtain $\|v_{tt}^N(t)\|^2 \leq L_4$, where $L_{3}, L_{4} $
are constants independent of $t$ and $N$.

From \eqref{e3.15}, \eqref{e3.19} and \eqref{e3.20}, we conclude that
\begin{equation} \label{e3.21}
\begin{gathered}
(u^{N})_{N} \text{ and } (v^N)_{N} \text{ are bounded in }
L^\infty((0,T), H^2(\Omega) \cap H_0^1 (\Omega)),\\
(u^N_{t})_{N} \text{ and } (v^N_{t})_{N} \text{ are bounded in }
L^\infty((0,T), H^1_0(\Omega)), \\
(u^N_{tt})_{N} \text{ and } (v^N_{tt})_{N} \text{ are bounded in }
 L^\infty((0,T), L^2(\Omega)).
\end{gathered}
\end{equation}

\noindent\textbf{Step 3.}
From \eqref{e3.21}, there exist subsequences of $(u^{N})_{N}$ and $(v^{N})_{N}$
(still denoted by $(u^{N})_{N}$ and $(v^{N})_{N})$, and two functions
$u, v : \Omega \times [0, T) \to \mathbb{R}$, such that
\begin{equation}
\begin{gathered}
u^{N} \overset{\star}{\rightharpoonup} u \text{ and }
 v^{N} \overset{\star}{\rightharpoonup} v \quad \text{in }
 L^\infty((0,T), H^2(\Omega) \cap H_0^1 (\Omega)), \\
u_{t}^{N} \overset{\star}{\rightharpoonup} u_{t} \text{ and }
 v_{t}^{N} \overset{\star}{\rightharpoonup} v_{t} \text{ in } L^\infty((0,T),
 H^1_0(\Omega)), \\
 u_{tt}^{N} \overset{\star}{\rightharpoonup} u_{tt} \text{ and }
 v_{tt}^{N} \overset{\star}{\rightharpoonup} v_{tt} \text{ in }
 L^\infty((0,T), L^2(\Omega)).
\end{gathered}\label{e3.22}
\end{equation}
Thanks to convergences \eqref{e3.2} and \eqref{e3.22}, taking $N \to +\infty$
in \eqref{e3.1}, we obtain
\begin{equation}
\begin{gathered}
\int_\Omega u_{tt}\phi_j dx + \int_\Omega \nabla u .\nabla \phi_j dx
- \int_\Omega \int_0^t g(t-s) \nabla u(s) .\nabla \phi_j\, ds\, dx
+ \alpha \int_\Omega v \phi_j \, dx = 0,\\
\int_\Omega v_{tt} \phi_j dx + \int_\Omega \nabla v.\nabla \phi_j dx
+ \alpha \int_\Omega u \phi_j \, dx = 0,\quad 1\leq j\leq N.
\end{gathered} \label{e3.23}
\end{equation}
 Consequently, for every $ \phi \in H_0^1(\Omega)\cap H^2(\Omega)$,
\begin{gather*}
\int_\Omega \Big[ u_{tt}- \Delta u + \int_0^t g(t-s) \Delta u(s)\, ds\,dx
 + \alpha v\Big] \phi \, dx = 0,\\
\int_\Omega [ v_{tt} - \Delta v + \alpha u] \phi \, dx = 0.
\end{gather*}
Therefore, $(u,v)$ solves the equations in \eqref{e1.1}.
\smallskip

\noindent\textbf{Step 4. Initial conditions.}
First, by Lions' lemma \cite{Lions} and since
\begin{gather*}
u^{N} \overset{\star}{\rightharpoonup} u \quad \text{in }
 L^\infty((0,T), H^2(\Omega) \cap H_0^1 (\Omega)), \\
u_{t}^{N} \overset{\star}{\rightharpoonup} u_{t} \quad \text{in }
 L^\infty((0,T), H^1_0(\Omega)),
\end{gather*}
we deduce, up to a subsequence, that
\begin{equation}
u^{N} \to u \quad \text{in } C([0, T], H^1_0(\Omega)).\label{e3.24}
\end{equation}
Therefore, \( u^{N}(\cdot, 0) \) is defined and
\[
u^{N}(\cdot, 0) \to u(\cdot, 0) \quad \text{in } H^1_0(\Omega).
\]
But,
\[ u^{N}(\cdot, 0) = u^{N}_0 \to u_0 \,\, \text{in}\,\, H^2(\Omega) \cap H^1_0(\Omega).\] Then \( u(\cdot, 0) = u_0 \, \). Similarly, we obtain $v(\cdot, 0)=v_0. $

Noting that
\[
 u_{tt}^{N} \overset{\star}{\rightharpoonup} u_{tt} \quad \text{and} \quad
 v_{tt}^{N} \overset{\star}{\rightharpoonup} v_{tt} \quad \text{in }
 L^\infty((0,T), L^2(\Omega)),
\]
we obtain $ u_t , v_{t} \in C([0, T), L^2(\Omega))$, as in lions \cite{Lions},
and thus \( u_t(\cdot, 0) \) has a meaning. Furthermore
\[
u^N_t(\cdot, 0) \to u_t(\cdot, 0) \quad \text{and} \quad
v^N_t(\cdot, 0) \to v_t(\cdot, 0) \quad\text{in } L^2(\Omega).
\]
But,
\[ u^N_t(\cdot, 0) = u^N_1 \to u_1 \quad \text{and} \quad
 v^N_t(\cdot, 0) = v^N_1 \to v_1 \quad \text{in } H^1_0(\Omega).
\]
Hence,
$u_t(x, 0) = u_1(x)$ and $v_t(x, 0) = v_1(x)$.
\smallskip

\noindent\textbf{Uniqueness.}
Assume that problem \eqref{e1.1} has two weak solutions, \( (u, v) \) and
\( (\overline{u},\overline{v}) \). Then,
\( (\mathbf{u}, \mathbf{v}) = (u - \overline{u}, v - \overline{v}) \) satisfies
\begin{gather}
 \mathbf{u}_{tt} - \Delta \mathbf{u} + \int_0^t g(t-s) \Delta \mathbf{u} \, ds
 + \alpha \mathbf{v} = 0, \quad \text{in } \Omega \times (0, T), \label{e3.25}\\
 \mathbf{v}_{tt} - \Delta \mathbf{v} + \alpha \mathbf{u} = 0, \quad
 \text{in } \Omega \times (0, T), \label{e3.26} \\
 \mathbf{u}(x,t) = \mathbf{v}(x,t) = 0, \quad \text{on } \partial\Omega \times (0, T),
 \label{e3.27}\\
 \mathbf{u}(x,0) = \mathbf{u}_t(x,0) =0,\, \mathbf{v}(x,0)=\mathbf{v}_t(x,0) = 0,
 \quad \text{in } \Omega. \label{e3.28}
\end{gather}
Multiplying \eqref{e3.25} by $\mathbf{u}_{t}$ and \eqref{e3.26} by $\mathbf{v}_{t}$,
integrating over $\Omega$, and adding the two resulting equations, we find
\begin{equation}
\begin{aligned}
&\frac{1}{2}\frac{d}{dt} \Big( \| \mathbf{u}_{t}(t) \| ^2
 + \Big[ 1-\int_0^t g(s) ds \Big] \| \nabla \mathbf{u} (t)\|^2
 +(g\circ \nabla \mathbf{u})(t) \\
&+2 \alpha \int_\Omega \mathbf{v}\textbf{ u}\,dx
 +\| \mathbf{v}_{t}(t) \|^2 + \| \nabla \mathbf{v} (t)\|^2 \Big)\\
& = -\frac{1}{2} g(t)\parallel \nabla \mathbf{u}(t) \parallel^2
 + \frac{1}{2}( g'\circ\nabla \mathbf{u}) (t).
\end{aligned} \label{e3.29}
\end{equation}
Using assumptions (A1)--(A3), and integrating over $(0, t),\, t\leq T$, we obtain
\[
 \| \mathbf{u}_{t}(t) \| ^2+ \Big[ 1-\int_0^t g(s) ds \Big]
 \| \nabla \mathbf{u}(t) \|^2 +(g\circ \nabla \mathbf{u})(t)
 +2 \alpha \int_\Omega \mathbf{v}\textbf{ u}dx
 +\| \mathbf{v}_{t} (t)\|^2 + \| \nabla \mathbf{v}(t) \|^2 \leq 0.
\]
 From Young and Poincar\'e's inequalities, we obtain
\begin{align*}
&\| \mathbf{u}_{t}(t) \| ^2+ \Big[ 1-\int_0^t g(s) ds
 -\alpha \varepsilon_{1} C_p \Big] \| \nabla \mathbf{u}(t) \|^2
 +(g\circ \nabla \mathbf{u})(t)\\
&+\| \mathbf{v}_{t}(t) \|^2 + \Big( 1-\frac{\alpha C_{p}}{\varepsilon_{1} }\Big)
\| \nabla \mathbf{v}(t) \|^2 \leq 0.
\end{align*}
We select $\varepsilon_{1}$ such that it satisfies\,
$\alpha C_p \leq \varepsilon_{1}\leq \frac{l}{\alpha C_p}$. Then
\[
 \| \mathbf{u}_{t} \|^2 + \| \nabla \mathbf{u} \|^2
 +(g\circ\nabla \mathbf{u})(t)+\| \mathbf{v}_{t} \|^2 + \| \nabla \mathbf{v} \|^2=0.
\]
Hence, \( \mathbf{u}_t(\cdot, t) = \mathbf{v}_t(\cdot, t) = 0 \) and
\( \nabla \mathbf{u}(\cdot, t) = \nabla\mathbf{v}(\cdot, t) = 0 \), for all
\( t \in (0, T) \). This implies that \( \textbf{ u} = \textbf{ v} = 0 \) on
\( \Omega \times (0, T) \), given that \( \mathbf{u} = \mathbf{v} = 0 \) on
\( \partial\Omega \times (0, T) \). This establishes the uniqueness.
\end{proof}

% Source-native exact authored range 634--1174
\section{Asymptotic behavior}

Now, we introduce the first and the second order energy functionals corresponding
to system \eqref{e1.1},
\begin{equation}
\begin{aligned}
E(t) &= \frac{1}{2} \| u_t(t) \|^2 + \frac{1}{2} \| v_t(t) \|^2
 + \frac{1}{2} \Big(1 - \int_0^t g(s) \, ds \Big) \| \nabla u(t) \|^2
 + \frac{1}{2} \| \nabla v(t) \|^2 \\
&\quad + \frac{1}{2} (g \circ \nabla u)(t) + \alpha \int_{\Omega} u v \, dx\label{e4.1}
\end{aligned}
\end{equation}
and
\begin{equation}
\begin{aligned}
 \mathsf{e}(t)
& = \frac{1}{2}\| \nabla u_t(t) \|^2+\frac{1}{2}\| \nabla v_t(t) \|^2
 +\frac{1}{2} \| \Delta v(t) \|^2
 + \frac{1}{2}\Big(1 - \int_0^t g(s) \, ds \Big) \| \Delta u(t) \|^2\\
&\quad + \frac{1}{2}(g \circ \Delta u)(t)
 + \alpha \int_{\Omega}\nabla v \nabla u dx.
\end{aligned}\label{e4.2}
\end{equation}

\begin{lemma} \label{lem4}
Under the conditions of Theorem 3, there exists a constant $m > 0$
such that, for all $ t \geq 0$,
\begin{gather}
E(t) \geq m \left( \| u_t(t) \|^2 + \| v_t(t) \|^2 + \| \nabla u(t) \|^2
+ \| \nabla v(t) \|^2 + (g \circ \nabla u)(t)\right), \label{e4.3} \\
\mathsf{e}(t) \geq m \left( \| \nabla u_t(t) \|^2+\| \nabla v_t(t) \|^2
+ \| \Delta v(t) \|^2 + \| \Delta u(t) \|^2
+ (g \circ \Delta u)(t)\right) .\label{e4.4}
\end{gather}
\end{lemma}

\begin{proof}
Applying Cauchy-Schwarz and Young's inequalities, for
$ \varepsilon_{1} > 0$, we obtain
\[
E(t) \geq \frac{1}{2} \Big( \| u_t(t) \|^2
 + \| v_t(t) \|^2 + \left[ l -\alpha \varepsilon_{1} C_p \right] \| \nabla u(t) \|^2
 + \big( 1-\frac{\alpha C_p}{\varepsilon_{1}}\big) \| \nabla v(t) \|^2
 + (g \circ \nabla u)(t)\Big),
 \]
 for all $t\geq 0$.
Keeping in mind \eqref{e2.3}, we can select $\varepsilon_{1}$ such that
 $\alpha C_p \leq \varepsilon_{1}\leq \frac{l}{\alpha C_p}$.
Consequently, we derive \eqref{e4.3} with
$ m = \frac{1}{2} \min \{ l -\alpha \varepsilon_{1} C_p,
1-\frac{\alpha C_p }{\varepsilon_{1}} \} $.
Using similar reasoning, \eqref{e4.4} can also be established with the same constant.
\end{proof}

\begin{lemma} \label{lem5}
Under the conditions of Theorem \ref{thm3}, the energy functional satisfies
\begin{gather}
E'(t) = \frac{1}{2} (g' \circ \nabla u)(t) - \frac{1}{2} g(t) \| \nabla u(t) \|^2
 \leq 0, \label{e4.5}\\
\mathsf{e}'(t)= \frac{1}{2}(g' \circ \Delta u)(t)-\frac{1}{2} g(t) \| \Delta u(t) \|^2
\leq 0.\label{e4.6}
\end{gather}
\end{lemma}

\begin{proof}
To obtain \eqref{e4.5}, we multiply the first equation of system \eqref{e1.1} by
$u_{t}$, the second one by $v_{t}$, we integrate each one over $\Omega$, utilizing
the boundary conditions, it follows that
\begin{align*}
&\frac{d}{dt} \Big( \frac{1}{2} \| u_{t}(t) \|^2 +\frac{1}{2} \| v_{t}(t) \|^2
+\frac{1}{2} \| \nabla u(t) \|^2 +\frac{1}{2} \| \nabla v(t) \|^2
 + \alpha \int_{\Omega} u\,v \Big) dx\\
&- \int_0^t g(t-s) \int_{\Omega} \nabla u(s) \nabla u_{t} \, dx ds =0.
\end{align*}
Applying Lemma \ref{lem1}, we arrive at
 \begin{align*}
&\frac{d}{2dt} \Big( \| u_{t}(t) \|^2 +\| v_{t}(t) \|^2
 + \Big[ 1-\int_0^t g(s) ds \Big] \| \nabla u(t) \|^2
 +\frac{1}{2} \| \nabla v(t) \|^2+ (g\circ \nabla u)(t)
 +\alpha \int_{\Omega} v \, u dx \Big) \\
& +\frac{1}{2} g(t)\parallel \nabla u(t) \parallel^2
- \frac{1}{2}( g'\circ\nabla u) (t) =0.
 \end{align*}
 To obtain \eqref{e4.6}, we multiply the equation \eqref{e1.1}$_1$ and
 \eqref{e1.1}$_2$ by $-\Delta u_t$ and $-\Delta v_t$, respectively,
and sum, to obtain
 \begin{align*}
&\frac{d}{dt}  \Big( \frac{1}{2} \| \nabla u_{t}(t) \|^2 + \frac{1}{2} \| \Delta u(t) \|^2 + \frac{1}{2}\| \nabla v_{t}(t) \|^2+ \frac{1}{2} \| \Delta v(t) \|^2 + \alpha \int_{\Omega} \nabla v \nabla u \Big) dx \\ 
& - \int_0^t g(t-s) \int_{\Omega} \Delta u(s) \Delta u_{t}(t) \, dx ds =0.
 \end{align*}
 Applying Lemma 1, we arrive at
\begin{align*}
&\frac{d}{2dt} \Big( \| \nabla u_t(t) \|^2+\| \nabla v_t(t) \|^2
+ \| \Delta v(t) \|^2 +\Big[ 1- \int_0^t g(s) ds \Big] \| \Delta u(t) \|^2 \\
& + (g \circ \Delta u)(t)+ 2 \alpha \int_{\Omega}\nabla v \nabla u dx\Big)
-\frac{1}{2}(g' \circ \Delta u)(t)+\frac{1}{2} g(t) \| \Delta u(t) \|^2=0.
\end{align*}
Then our conclusion follows.
\end{proof}

Now we construct a Lyapunov functional $L$ equivalent to $E+e$, with which we can
show the desired result.
Similarly to \cite{JKJT}, we define
\[
 C_\mu = \int_0^{\infty} \frac{g^2(s)}{\mu g(s) - g'(s)} \, ds, \quad
 h(t) = \mu g(t) - g'(t)\quad
 \text{for } 0<\mu <1.
\]

\begin{remark}[\cite{MustafaMuhammadR}] \label{rmk6} \rm
(1) The well-known Jensen's inequality will be of essential use in establishing our
 main result. For the sake of completeness, let us present it here.
 If \( F \) is a convex function on \( [a, b] \), \( f : \Omega \to [a, b] \), and
 \( h \) are integrable functions on \( \Omega \) with
 \( h(x) \geq 0 \) and $\int_\Omega h(x)\,dx = \kappa > 0$,
then Jensen's inequality states that
\[
F\Big( \frac{1}{\kappa} \int_\Omega f(x) h(x)\,dx \Big)
\leq \frac{1}{\kappa} \int_\Omega F(f(x)) h(x)\,dx.
\]

(2) From condition (A1), it follows that $ \lim_{t \to +\infty} g(t) = 0$.
Thus, there exists \( t_0 \geq 0 \) large enough such that
\[
g(t_0) = r \quad \text{and} \quad g(t) \leq r, \, \forall t \geq t_0.
\]
Since $g$ and $\zeta$ are positive, non-increasing and continuous functions, and
$G$ is a positive and continuous function, we have for all \( t \in [0, t_0] \),
\[
0 < g(t_0) \leq g(t) \leq g(0), \quad 0 < \zeta(t_0) \leq \zeta(t) \leq \zeta(0).
\]
This implies the existence of two positive constants \( a \) and \( b \) such that
\[
a \leq \zeta(t) G(g(t)) \leq b.
\]
As a result, we have, for all \( t \in [0, t_0] \),
\[
g'(t) \leq -\zeta(t) G(g(t)) \leq -a,
\]
consequently,
\[
g'(t) \leq -\frac{a}{g(0)} g(0) \leq -\frac{a}{g(0)} g(t). \label{e4.7}
\]

(3) Given that
\[
\frac{\mu g^2(s)}{\mu g(s) - g'(s)} < g(s),
\]
and applying the Lebesgue dominated convergence theorem, we easily check that
\[
\mu C_\mu = \int_0^\infty \frac{\mu g^2(s)}{\mu g(s) - g'(s)} \, ds \to 0
\quad \text{as } \mu \to 0.
\]

(4) If $G$ is a strictly increasing and strictly convex $C^2$ function on $(0, r]$,
with $G(0) = G'(0) = 0$, then it has an extension $\overline{G}$ which is strictly
increasing and strictly convex $C^2$ function on $(0, \infty)$.
See \cite{MustafaMuhammadR}.
\end{remark}

\begin{lemma} \label{lem7}
Let $(u,v)$ be the solution of \eqref{e1.1}. Consider the functional
\[
 \psi(t) = \int_{\Omega} uu_t \, dx +\int_{\Omega} vv_t \, dx.
\]
Then, for any \(\epsilon > 0\),
\begin{equation}
 \psi '(t) \leq \| u_t(t) \|^2 +\| v_t(t) \|^2 - (l-\epsilon) \| \nabla u(t) \|^2 - \| \nabla v (t)\|^2 -2 \alpha \int_{\Omega} u \, v dx + \frac{ C_\mu }{4\epsilon} (h \circ \nabla u)(t).
 \label{e4.8}
\end{equation}
\end{lemma}

\begin{proof}
Using system \eqref{e1.1}, a direct computation leads to
\begin{equation}
\begin{aligned}
 \psi'(t) &= \| u_t(t) \|^2 + \| v_t(t) \|^2 -\| \nabla u(t) \|^2
  -\| \nabla v(t) \|^2 + \int_{\Omega} \nabla u(t)
  \Big( \int_0^t g(t - s) \nabla u(s) ds \Big) dx \\
&\quad - 2 \alpha\int_{\Omega} u \, v dx.
\end{aligned} \label{e4.9}
\end{equation}
Using Young's inequality and the fact that
$\int_0^t g(s) \, ds \leq \int_0^{\infty} g(s) \, ds = 1 - l$, we
estimate the fifth term in \eqref{e4.9} to  obtain
\begin{equation}
\begin{aligned}
 \psi'(t) &\leq \| u_t(t) \|^2 + \| v_t(t) \|^2 - l \| \nabla u(t)\|^2
 -\| \nabla v(t) \|^2 + \epsilon \| \nabla u(t) \|^2\\
&\quad +\frac{1}{4\epsilon} \int_{\Omega}
 \Big( \int_0^t g(t - s) |\nabla u(s) - \nabla u(t)| ds \Big)^2 dx
 - 2 \alpha\int_{\Omega} u \, v dx.
\end{aligned}\label{e4.10}
\end{equation}
 Now, the  Cauchy-Schwarz inequality gives
\begin{equation}
\begin{aligned}
&\int_{\Omega} \Big( \int_0^{t} g(t - s) |\nabla u(s) - \nabla u(t)| \, ds \Big)^2
 \, dx \\
&=\int_{\Omega} \Big(\int_0^{t} \frac{g(t - s)}{\sqrt{\mu g(t - s)
 - g'(t - s)}} \sqrt{\mu g(t - s) - g'(t - s)} |\nabla u(s) - \nabla u(t)| \, ds \Big)^2
   \, dx\\
&\leq \Big( \int_0^{t} \frac{g^2(s)}{\mu g(s) - g'(s)} \, ds \Big)
 \int_{\Omega} \Big( \int_0^{t} [\mu g(t - s) - g'(t - s)] |\nabla u(s)
 - \nabla u(t)|^2 \, ds \Big) \, dx \\
&\leq C_\mu (h \circ \nabla u)(t).
\end{aligned} \label{e4.11}
\end{equation}
By substituting \eqref{e4.11} into \eqref{e4.10}, we obtain \eqref{e4.8}.
\end{proof}

\begin{lemma}  \label{lem8}
Let $(u, v)$ be the solution of \eqref{e1.1}. Then, the functional
\[
 \chi(t) = -\int_{\Omega} u_t \int_0^t g(t - s)(u(t) - u(s)) \, ds \, dx,
\]
satisfies, for any $ \delta_{1} > 0$, the estimate
\begin{equation}
\begin{aligned}
 \chi'(t) &\leq  \Big(\delta_{1} - \int_0^t g(s) \, ds \Big) \| u_t (t)\|^2
  + \delta_{1} \| \nabla u (t) \|^2
  +\alpha^2 C_p \frac{\delta_{1}}{2} \| \nabla v (t)\|^2 \\
&\quad +\Big(C+ C_\mu \Big( 1+ \frac{C}{ \delta_{1}}\Big) \Big)
 (h \circ \nabla u)(t),
\end{aligned} \label{e4.12}
\end{equation}
 where $C=\max\{c_1,c_2\},$\quad with $c_1=\frac{C_{p}\big(\alpha(1-l)+g(0)\big)}{2}$,
 $c_2= c+\frac{C_{p}(1+\mu^2)}{2}$.
\end{lemma}

\begin{proof}
From \eqref{e1.1} and integration by parts, we obtain
\begin{equation}
\begin{aligned}
 \chi'(t) &= \Big(1 - \int_0^t g(s) \, ds\Big)
 \int_{\Omega} \nabla u(t) \, \int_0^t g(t - s)(\nabla u(t) - \nabla u(s)) \, ds \, dx \\
&\quad + \int_{\Omega} \Big( \int_0^t g(t - s)(\nabla u(s)
 - \nabla u(t)) \, ds\Big)^2 \, dx
 - \int_{\Omega} u_t \, \int_0^t g'(t - s)(u(t) - u(s)) \, ds \, dx \\
&\quad - \Big( \int_0^t g(s) \, ds \Big) \| u_t(t) \|^2
 + \alpha \int_{\Omega} v \, \int_0^t g(t - s)(u(t) - u(s)) \, ds \, dx.
\end{aligned} \label{e4.13}
\end{equation}
 By Cauchy-Schwartz inequality, the second term of \eqref{e4.13} satisfies
\begin{equation}
\int_{\Omega} \Big( \int_0^t g(t - s)(\nabla u(s)
 - \nabla u(t)) \, ds \Big)^2 \, dx
\leq C_\mu (h \circ \nabla u)(t). \label{e4.14}
\end{equation}
Similarly, for any $\delta_{1} > 0$, we have
\begin{equation}
\begin{aligned}
&\Big(1 - \int_0^t g(s) \, ds \Big) \int_{\Omega} \nabla u
 \int_0^t g(t - s)(\nabla u(t) -\nabla u(s)) \, ds \, dx \\
&\leq \delta_{1} \| \nabla u(t) \|^2
 + \frac{\left(1 - g_0\right)^2 C_\mu}{4\delta_{1}} \, (h \circ \nabla u)(t)\\
&\leq \delta_{1} \| \nabla u(t) \|^2 + \frac{c C_\mu}{\delta_{1}}
 (h \circ \nabla u)(t),
\end{aligned} \label{e4.15}
\end{equation}
\begin{equation}
\begin{aligned}
& -\int_{\Omega} u_t \, \int_0^t g'(t - s)(u(t) - u(s)) \, ds \, dx\\
&= \int_{\Omega} u_t \, \int_0^t h(t - s)(u(t) - u(s)) \, ds \, dx
 - \int_{\Omega} u_t \, \int_0^t \mu g(t - s)(u(t) - u(s)) \, ds \, dx \\
&\leq \delta_{1} \| u_t(t) \|^2+ \frac{1}{2\delta_{1}} \int_{\Omega}
\Big( \int_0^{t} \sqrt{h(t - s)} \sqrt{h(t - s)} |u(s) - u(t)| \, ds \Big)^2 \, dx \\
&\quad + \frac{\mu^2}{2\delta_{1}} \int_{\Omega}
 \Big( \int_0^{t} g(t - s) |u(s) - u(t)| \, ds \Big)^2 \, dx \\
&\leq \delta_{1} \| u_t(t) \|^2+ \frac{ \int_0^{t} h(s) \, ds }{2\delta_{1}}
 (h \circ u)(t)+ \frac{\mu^2 C_\mu }{2\delta_{1}} (h \circ u)(t)\\
&\leq \delta_{1} \| u_t(t) \|^2+ \frac{C_p \big(\alpha(1-l)+g(0)\big)}{2\delta_{1}}
 (h \circ \nabla u)(t)+ \frac{C_p\mu^2 C_\mu }{2\delta_{1}} (h \circ \nabla u)(t)
\end{aligned} \label{e4.16}
\end{equation}
and
\begin{equation}
 \alpha \int_{\Omega} v  \int_0^t g(t - s)(u(t) - u(s)) \, ds \, dx
 \leq \alpha^2 C_p \,\frac{\delta_{1}}{2} \| \nabla v(t) \|^2
 + C_p \frac{ C_\mu}{ 2 \delta_{1}} (h \circ \nabla u)(t).\label{e4.17}
\end{equation}

Collecting  \eqref{e4.14}--\eqref{e4.17}, we obtain the required estimate.
\end{proof}

\begin{lemma} \label{lem9}
Assume that {\rm (A1)--(A3)} hold. Then, the functional
\begin{equation*}
 I(t) = \int_{\Omega} u_{tt} \, v_t dx - \int_{\Omega} v_{tt} \, u_t dx
 \end{equation*}
 satisfies along the solution of \eqref{e1.1} and for any $\delta_2>0$,
 \begin{equation}
 I'(t) \leq - ( \alpha- \delta_2) \| v_t (t) \|^2+ \alpha \| u_t (t) \|^2
 + \frac{C g^2(t)}{\delta_2} \| \Delta u(t) \|^2-\frac{C}{ \delta_2}
 (g' \circ \Delta u)(t). \label{e4.18}
\end{equation}
\end{lemma}

\begin{proof}
By using equations \eqref{e1.1} and integrating by parts, we have
\begin{align*}
 \frac{d}{dt}\int_{\Omega} u_{tt} \, v_t dx
& =\int_{\Omega} u_{tt} \, v_{tt} dx -\int_{\Omega} \nabla u_t \nabla v_t dx
 - \alpha \| v_t (t) \|^2 \\
&\quad  +\int_{\Omega} v_t \int_0^t g'(t - s)(\Delta u(t) - \Delta u(s)) \, ds\,dx
- g(t) \int_{\Omega} \Delta u v_t dx,
\end{align*}
and
\begin{equation*}
- \frac{d}{dt}\int_{\Omega} v_{tt} \, u_t dx= -\int_{\Omega} v_{tt} \, u_{tt} dx
+\int_{\Omega} \nabla u_t \nabla v_t dx + \alpha \| u_t (t)\|^2.
\end{equation*}
Addition of the last two identities leads to
\[
I'(t)= \int_{\Omega} v_t \int_0^t g'(t - s)(\Delta u(t)
- \Delta u(s)) \, ds\,dx- g(t) \int_{\Omega} \Delta u v_t dx
 - \alpha \| v_t(t) \|^2 + \alpha \| u_t(t) \|^2
\]
and Young's inequality yields
\begin{gather*}
\int_{\Omega} v_t \int_0^t g'(t - s)(\Delta u(t) - \Delta u(s)) \, ds\,dx
\leq \frac{\delta_2}{2} \| v_t(t) \|^2- \frac{ C}{\delta_2}
(g' \circ \Delta u)(t),
\\
 g(t) \int_{\Omega} \Delta u v_t dx \leq \frac{\delta_2}{2} \| v_t(t) \|^2
 + \frac{Cg^2(t)}{ \delta_2} \| \Delta u (t) \|^2.
\end{gather*}
 We finally get the estimate \eqref{e4.18}.
\end{proof}

\begin{lemma}[\cite{MustafaMuhammadR}] \label{lem10}
 Assume that {\rm (A1)--(A3)}  hold.
Then
\[
 K(t) = \int_{\Omega} \int_0^t f(t - s) |\nabla u(s)|^2 \, ds \, dx,
\]
where $f(t) = \int_{t}^{\infty} g(s) \, ds$, satisfies along the solution
of \eqref{e1.1}, the estimate
\begin{equation}
K'(t) \leq -\frac{1}{2} (g \circ \nabla u)(t)
+ 3(1 - l) \int_{\Omega} |\nabla u(t)|^2 \, dx.   \label{e4.19}
\end{equation}
\end{lemma}

At this stage, we introduce the Lyapunov functional
\begin{equation}
 L(t) = \eta( E(t)+ e(t) )+ \eta_1 \psi(t) + \eta_2 \chi(t)+ I(t), \label{e4.20}
\end{equation}
where $ \eta , \eta_1, \eta_2> 0$, are constants to be suitably specified later.

\begin{proposition} \label{prop11}
There exist \(\rho_i> 0 \), (\( i = 1, 2 \)) such that
\begin{equation}
\rho_1 (E(t)+e(t)) \leq L(t) \leq \rho_2 (E(t)+e(t)),\, t\geq 0. \label{e4.21}
\end{equation}
\end{proposition}

\begin{proof}
Using Young's and Cauchy-Schwarz inequalities, we see that
\begin{align*}
\psi(t)
& \leq \frac{1}{2}\int_{\Omega} |u|^2 \, dx+ \frac{1}{2}\int_{\Omega}|u_t|^2 \, dx
+\frac{1}{2}\int_{\Omega} |v|^2 \, dx+ \frac{1}{2}\int_{\Omega}|v_t|^2 \, dx \\
&\leq \frac{C_p}{2} \| \nabla u(t) \|^2+ \frac{1}{2} \| u_t(t) \|^2
 + \frac{C_p}{2} \| \nabla v(t) \|^2+ \frac{1}{2} \| v_t(t) \|^2,
\end{align*}
and
\begin{align*}
 \chi(t)
&\leq \frac{1}{2}\int_{\Omega}|u_t|^2 \, dx
 +\frac{1}{2} \int_{\Omega} \Big( \int_0^{t} g(t - s)|u(t) - u(s)| \, ds \Big)^2 \, dx
 \\
&\leq \frac{1}{2} \| u_t(t) \|^2+\frac{C_p^2 }{2} (1-l) (g \circ \nabla u)(t).
\end{align*}
Similarly, using \eqref{e1.1}, we obtain
\begin{align*}
I(t)
& \leq \frac{1}{2} \| \Delta u(t) \|^2+\frac{1}{2}\| v_t (t)\|^2
 +\int_{\Omega} | v_t| \int_0^t g(t-s) |\Delta u(s)- \Delta u(t)|\, ds\, dx
 +\frac{(1-l)}{2} \| v_t(t) \|^2\\
&\quad +\frac{(1-l)}{2} \| \Delta u (t)\|^2
 +\frac{\alpha}{2}C_p \| \nabla v(t) \|^2+\frac{\alpha}{2}\| v_t(t) \|^2
  +\frac{1}{2}\| \Delta v(t) \|^2+\frac{1}{2}\| u_t(t) \|^2\\
&\quad  +\frac{\alpha}{2}C_p \| \nabla u(t) \|^2+\frac{\alpha}{2} \| u_t(t) \|^2\\
& \leq \Big( \frac{1}{2}+\frac{(1-l)}{2}\Big)
 \| \Delta u(t)\|^2+\Big(1+\frac{(1-l)}{2}+\frac{\alpha}{2}\Big)
  \| v_t(t) \|^2 +\frac{(1-l)}{2} (g \circ \Delta u)(t)\\
&\quad + \frac{\alpha}{2}C_p \| \nabla v (t)\|^2
 +\frac{\alpha}{2}C_p \| \nabla u (t)\|^2+\frac{1}{2}\| \Delta v(t) \|^2
 +\Big( \frac{1}{2}+\frac{\alpha}{2} \Big) \| u_t (t)\|^2.
\end{align*}
Then
\begin{align*}
|L(t) - \eta (E(t)+ e(t))|
& \leq \frac{C_p }{2}( \eta_1 +\alpha )\| \nabla u(t) \|^2
 +\Big( \frac{\eta_1}{2} +\frac{\eta_2}{2}+\frac{\alpha}{2}+\frac{1}{2}\Big)
  \| u_t(t) \|^2+ \frac{C_p \eta_1}{2} \| \nabla v (t)\|^2\\
&\quad +\Big( \frac{\eta_1}{2}+1+\frac{(1-l)}{2} +\frac{\alpha}{2}\Big)
 \| v_t(t) \|^2+\frac{ C_p^2 \eta_2}{2} (1-l) (g \circ \nabla u)(t)\\
&\quad  +\Big( \frac{1}{2}+\frac{(1-l)}{2} \Big) \| \Delta u(t)\|^2
 +\frac{1}{2}\| \Delta v(t) \|^2 + \frac{(1-l)}{2} (g \circ \Delta u)(t)\\
&\leq C \left(\eta_1+\eta_2\right)\Big( \| \nabla u(t) \|^2+\| u_t(t) \|^2
 +\| \nabla v(t) \|^2+\| v_t(t) \|^2+ (g \circ \nabla u)(t)\\
&\quad +\| \Delta u(t)\|^2+\| \Delta v(t) \|^2+(g \circ \Delta u)(t)\Big).
\end{align*}
Therefore,
\[
|L(t) - \eta (E(t)+ e(t))|  \leq C\left(\eta_1+\eta_2\right)(E(t)+ e(t)),
\]
which yields
\[
(\eta-C\left(\eta_1+\eta_2\right))(E(t)+ e(t))
\leq L(t)\leq (\eta+C\left(\eta_1+\eta_2\right))(E(t)+ e(t)).
\]
So, by choosing $\eta > C \left(\eta_1+\eta_2\right)$\,, we obtain \eqref{e4.21}.
\end{proof}


\begin{lemma} \label{lem12}
Assume that conditions {\rm (A1)--(A3)} hold. Then, for any $ M, N>0$ and for
suitable choice of $\eta, \eta_1, \eta_2 >0$,
we have
\begin{equation}
\begin{aligned}
 L'(t)
&\leq  - \frac{l\alpha}{8} \| v_t (t) \|^2
 - \Big( \frac{3 \alpha}{2}- \frac{3\alpha l}{N}\Big) \| u_t(t) \|^2
 - \alpha^2 \int_{\Omega} u \, v dx
 - \Big( \frac{\alpha}{2} (l- \epsilon)- \frac{3\alpha l}{N} \Big)
 \| \nabla u(t) \|^2\\
&\quad - \Big( \frac{\alpha}{2} - 3 C_p \alpha^{3} \frac{l}{2N} \Big)
 \| \nabla v (t)\|^2+\frac{1}{2M} (g \circ \nabla u)(t), \quad \text{for all }
  t \geq t_0,
\end{aligned}\label{e4.22}
\end{equation}
where $t_0$ is introduced in remark \ref{rmk6}.
\end{lemma}

\begin{proof}
By combining \eqref{e4.5}, \eqref{e4.6}, \eqref{e4.8}, \eqref{e4.12} and
\eqref{e4.18}, we obtain
\begin{align*}
L'(t)
&\leq \frac{\eta}{2} (g' \circ \nabla u)(t) - \frac{\eta}{2} g(t) \| \nabla u(t) \|^2
 + \frac{\eta}{2}(g' \circ \Delta u)(t)-\frac{\eta}{2} g(t) \| \Delta u(t) \|^2
 + \eta_1 \| u_t (t)\|^2 + \eta_1 \| v_t(t) \|^2 \\
&\quad -\eta_1\| \nabla v(t) \|^2 - \eta_1 (l-\epsilon) \| \nabla u(t) \|^2
 -2 \alpha \eta_1 \int_{\Omega} u  v \,dx
 + \frac{ C_\mu \eta_1}{4\epsilon} (h \circ \nabla u)(t)
 + \eta_2 \delta_{1} \| \nabla u(t) \|^2\\
&\quad + \eta_2 \Big(\delta_{1} - \int_0^t g(s) \, ds \Big)\| u_t(t) \|^2
 +\eta_2 \alpha^2 C_p \frac{\delta_{1}}{2} \| \nabla v(t) \|^2
 + \eta_2 \Big(C+ C_\mu \Big( 1+ \frac{C}{ \delta_{1}}\Big) \Big)(h \circ \nabla u)(t)\\
&\quad - ( \alpha - \delta_2) \| v_t(t) \|^2+ \alpha \| u_t(t) \|_2^2
 + \frac{C g^2(t)}{\delta_2} \| \Delta u (t) \|_2^2-\frac{C}{ \delta_2}
 (g' \circ \Delta u)(t).
\end{align*}
Letting $g_0=\int_0^{t_0} g(s) \, ds$ and
recalling that $g'=(\mu g- h)$, we obtain
\begin{align*}
 L'(t)
&\leq  -( \alpha- \delta_2- \eta_1) \| v_t (t) \|^2
 - \left( \eta_2 \left(g_0- \delta_{1} \right)-\eta_1-\alpha \right)
  \| u_t(t) \|^2- 2 \alpha \eta_1 \int_{\Omega} u \, v dx \\
&\quad  - \left( \eta_1 (l- \epsilon)- \eta_2 \delta_{1} \right) \| \nabla u(t) \|^2
 - \Big( \eta_1- \eta_2 C_p \alpha^2 \frac{\delta_{1}}{2} \Big)
 \| \nabla v (t)\|^2+\Big(\frac{\eta}{2}-\frac{C}{ \delta_2} \Big)
  (g' \circ \Delta u)(t) \\
&\quad - g(t)\Big(\frac{\eta}{2}-\frac{Cg(t)}{\delta_2}\Big)
 \| \Delta u(t) \|^2- \Big( \frac{\eta}{2}
 - C \eta_2 - C_\mu \Big( \frac{\eta_1}{4\epsilon}+\eta_2
  \Big( 1+ \frac{C}{ \delta_{1}}\Big) \Big)\Big)(h \circ \nabla u)(t)\\
&\quad - \frac{\eta}{2} g(t) \| \nabla u (t)\|^2
 + \frac{\eta}{2} \mu (g \circ \nabla u)(t), \quad \forall t \geq t_0.
\end{align*}
Taking $\delta_2\leq \frac{\alpha l}{4}$,
 $\delta_{1}=\frac{l g_0}{N}$ ($N\geq 12$, to be chosen later) and
 $\eta_2=\frac{3 \alpha }{g_0}$, we find that for all $t\geq t_0$,
\begin{align*}
 L'(t)
&\leq  -( \alpha-\delta_2- \eta_1) \| v_t(t) \|^2
 - \Big( 2 \alpha- \frac{3\alpha l}{N} -\eta_1 \Big)
  \| u_t(t) \|^2-2 \alpha \eta_1 \int_{\Omega} u \, v dx \\
 &\quad - \Big( \eta_1 (l- \epsilon)- \frac{3\alpha l}{N} \Big) \| \nabla u(t) \|^2
 - \Big( \eta_1- 3 C_p \alpha^{3} \frac{l}{2N} \Big)
  \| \nabla v(t) \|^2+\Big(\frac{\eta}{2}-\frac{C}{\delta_2} \Big)
   (g' \circ \Delta u)(t) \\
 &\quad - g(t)\Big(\frac{\eta}{2}-\frac{Cg(t)}{\delta_2}\Big) \| \Delta u(t) \|^2
 - \Big( \frac{\eta}{2} - C \frac{3 \alpha }{g_0}
  - C_\mu \Big( \frac{\eta_1}{4\epsilon}+\frac{3 \alpha }{g_0}
  \Big( 1+ \frac{N C }{ l g_0}\Big) \Big)\Big)(h \circ \nabla u)(t) \\
&\quad +\frac{\eta\mu}{2} (g \circ \nabla u)(t)- \frac{\eta}{2} g(t) \| \nabla u(t) \|^2.
\end{align*}
For \(0< \epsilon < \frac{l^2}{2}\), we choose $\eta_{1}$ such that
\begin{equation}
\begin{gathered}
(l- \epsilon)\eta_1 - \frac{3\alpha l}{N} > \frac{l\alpha}{4}(1-l),\quad
\alpha - \delta_2 - \eta_1 > \frac{l\alpha}{8},\\
2\alpha-\frac{3\alpha l}{N} - \eta_1 > \frac{\alpha (1+l)}{4},\quad
\eta_1 -3 C_p \alpha^{3} \frac{l}{2N} > \frac{\alpha}{4}.
\end{gathered}\label{e4.23}
\end{equation}
In particular, we take $\eta_1=\frac{\alpha}{2}$ and \(N\) large so that
\[
\frac{\alpha}{2} - 3 C_p \alpha^{3} \frac{l}{2N} > \frac{\alpha}{4},
\]
and the inequalities \eqref{e4.23} are satisfied.
\\
 By applying Remark \ref{rmk6}, there exists \(0 < \mu_0 < 1\), such that
if \(\mu < \mu_0\), then
 \[
\mu C_\mu < \frac{1}{8\big( \frac{\alpha}{16\epsilon}+\frac{3 \alpha }{g_0}
\big( 1+ \frac{N C }{ l g_0}\big) \big)}.
\]
Now, we choose $\eta$ large enough and $\mu$ so that
\[
\frac{\eta}{2}-\frac{C}{\delta_2}> 0, \quad
\frac{\eta}{2} - \frac{3 \alpha C}{g_0} >0 \quad \text{and}\quad
\mu = \frac{1}{M\eta} < \mu_0,
\]
which gives
\[
 \frac{\eta}{2} - C \frac{3 \alpha }{g_0}
 - C_\mu \Big( \frac{\alpha}{16\epsilon}+\frac{3 \alpha }{g_0}
 \Big( 1+ \frac{N C}{ l g_0}\Big) \Big)>0.
\]
Hence,
\begin{align*}
 L'(t)
&\leq  - \frac{l\alpha}{8} \| v_t(t) \|^2
 - \Big( \frac{3 \alpha}{2}- \frac{3\alpha l}{N}\Big) \| u_t(t) \|^2
 - \alpha^2 \int_{\Omega} u  v \,dx
 - \Big( \frac{\alpha}{2} (l- \epsilon)- \frac{3\alpha l}{N} \Big)
  \| \nabla u(t) \|^2\\
&\quad - \Big( \frac{\alpha}{2} - 3 C_p \alpha^{3} \frac{l}{2N} \Big)
\| \nabla v (t)\|^2+\frac{1}{2M} (g \circ \nabla u)(t).
\end{align*}
We further choose $\epsilon$ even smaller, $N$ and $\eta$ larger, such
that \eqref{e4.21} remains valid, we obtain \eqref{e4.22}.
\end{proof}

% Source-native exact authored range 1175--1186
\begin{theorem} \label{thm13}
Assume that {\rm (A1)--(A3)} hold and $u_0,v_0 \in H^2(\Omega) \cap H^1_0(\Omega)$,
$u_1,v_1 \in H^1_0(\Omega)$. Then, there exist positive constants $k_0$, $C$ and
$\gamma$ such that the energy functional $E$ satisfies
\begin{gather}
E(t) \leq k_0 \frac{\left( E + e \right)(t_0)}{\int_{t_0}^t\xi(s)ds},\quad
\forall t\geq t_0, \quad  \text{if $G$ is linear.} \label{e4.24}\\
E(t)\leq C G_2^{-1}\Big( \frac{\gamma}{\int_{t_0}^t \xi(s)ds}\Big), \quad
 \forall t>t_0, \quad \text{if $G$ is non linear}, \label{e4.25}
\end{gather}
where $ G_2(\tau ) = \tau G'(\tau )$.
\end{theorem}
% Source-native exact authored range 1188--1431
\begin{proof}
Recalling \eqref{e4.22}, then for some constant $k > 0$ and for all $t \geq t_0$,
we obtain
\begin{equation}
\begin{aligned}
L'(t) &\leq - \frac{l\alpha}{8} \| v_t(t) \|^2
 - \Big( \frac{3 \alpha}{2}- \frac{3\alpha l}{N}\Big) \| u_t(t) \|^2
  - \alpha^2 \int_{\Omega} u \, v dx
 - \Big( \frac{\alpha}{2} (l- \epsilon)- \frac{3\alpha l}{N} \Big) \| \nabla u(t) \|^2\\
&- \Big( \frac{\alpha}{2} - 3 C_p \alpha^{3} \frac{l}{2N} \Big)
 \| \nabla v (t)\|^2 +\frac{1}{2M} (g \circ \nabla u)(t)\\
& \leq -k E(t) + c (g \circ \nabla u)(t).
\end{aligned} \label{e4.26}
\end{equation}
\smallskip

\noindent\textbf{Case (I): $G(t)$ is linear.}
Multiplying \eqref{e4.26} by $\xi(t)$ and using (A1) and \eqref{e4.5},
we find
\begin{align*}
\xi(t) L'(t)
&\leq -k \xi(t) E(t) + c \xi(t) \int_0^{t} g(s)
 \int_{\Omega} |\nabla u(t) - \nabla u(t - s)|^2 \, dx \, ds\\
& \leq -k \xi(t) E (t) - c E'(t),\quad \forall t\geq t_0.
\end{align*}
As $\xi$ is non-increasing and $L \geq 0$, we reach
\[
(\xi(t)L(t))'\leq \xi(t)L'(t)\leq - k \xi(t) E(t) - c E'(t),\quad \forall t\geq t_0;
\]
which yields
\[
\xi(t) E(t)\leq -\frac{1}{k}(\xi L+c E)'(t),\quad \forall t\geq t_0.
\]
Simple integration over $(t_0, t)$ leads to
\begin{align*}
E(t) \int_{t_0}^t\xi(s)ds
&\leq \int_{t_0}^t \xi(s) E(s)ds \\
&\leq -\int_{t_0}^t \frac{1}{k}(\xi L+c E)'(s)ds \\
& \leq \frac{1}{k}(\xi L+c E)(t_0),\quad \forall t\geq t_0.
\end{align*}
Using $\xi L+c E \sim E+e$, estimate \eqref{e4.24} is established.
\smallskip

\noindent\textbf{Case (II): $G(t)$  nonlinear.}
As in \cite{MM181}, we use \eqref{e4.5} and \eqref{e4.7} to deduce that,
for any $t \geq t_0$,
\begin{align*}
&\int_0^{t_0} g(s) \int_{\Omega} |\nabla u(t) - \nabla u(t - s)|^2 \, dx \, ds\\
&\leq - \frac{g(0)}{a} \int_0^{t_0} g'(s) \int_{\Omega} |\nabla u(t)
- \nabla u(t - s)|^2 \, dx \, ds \\
&\leq -cE'(t),
\end{align*}
which can be utilized in \eqref{e4.26}, to obtain, for all $t \geq t_0$,
\begin{align*}
L'(t)
&\leq -k E(t) + c (g \circ \nabla u)(t)\\
& \leq -kE (t) - c E'(t)+ c \int_{t_0}^{t} g(s) \int_{\Omega} |\nabla u(t)
- \nabla u(t - s)|^2 \, dx \, ds.
\end{align*}
By taking $F (t) = L(t) + c E(t)$, which is equivalent to $E(t)+e(t)$, we infer that
\begin{equation}
F'(t) \leq -k E(t) + c \int_{t_0}^{t} g(s) \int_{\Omega} |\nabla u(t)
- \nabla u(t - s)|^2 \, dx \, ds.\label{e4.27}
\end{equation}
We use Lemmas \ref{lem10} and \ref{lem12} to conclude that, for $\beta>0$,
$$
\mathcal{L}(t) = L(t) +\beta K(t)
$$
is nonnegative and  for all $t \geq t_0$ satisfies
\begin{align*}
 \mathcal{L}'(t)
&\leq  - \frac{l\alpha}{8} \| v_t(t) \|^2
 - \Big( \frac{3 \alpha}{2}- \frac{3\alpha l}{N}\Big)
  \| u_t(t) \|^2- \alpha^2 \int_{\Omega} u \, v dx
 - \Big( \frac{\alpha}{2} (l- \epsilon)- \frac{3\alpha l}{N} \Big) \| \nabla u(t) \|^2\\
&\quad - \Big( \frac{\alpha}{2} - 3 C_p \alpha^{3} \frac{l}{2N} \Big)
  \| \nabla v(t) \|^2 +\frac{1}{2M} (g \circ \nabla u)(t) -\frac{\beta}{2}
  (g \circ \nabla u)(t)+3 \beta (1-l)\| \nabla u(t) \|^2\\
&\leq  - \frac{l\alpha}{8} \| v_t(t) \|^2
 - \Big( \frac{3 \alpha}{2}- \frac{3\alpha l}{N}\Big)
  \| u_t(t) \|^2- \alpha^2 \int_{\Omega} u \, v dx  \\
&\quad  - \Big( \frac{\alpha}{2} (l- \epsilon)- \frac{3\alpha l}{N}- 3 \beta (1-l) \Big)
  \| \nabla u(t) \|^2\\
&\quad - \Big( \frac{\alpha}{2} - 3 C_p \alpha^{3} \frac{l}{2N} \Big)
  \| \nabla v(t) \|^2
 - \Big(\frac{\beta}{2}-\frac{1}{2M}\Big)(g \circ \nabla u)(t).
\end{align*}
By selecting $\beta $ sufficiently small and then $M $ sufficiently large
such that
\begin{gather*}
\frac{\alpha}{2} (l- \epsilon)- \frac{3\alpha l}{N}- 3 \beta (1-l) >0,\\
\frac{\beta}{2}-\frac{1}{2M}>0,
\end{gather*}
we easily obtain
$\mathcal{L}'(t)\leq -d E(t)$,
where $d$ is some positive constant. So,
\[
d \int_{t_0}^t E(s) \, ds \leq \mathcal{L}(t_0) - \mathcal{L}(t)
\leq \mathcal{L}(t_0),
\]
which implies that
\begin{equation}
\int_0^{\infty} E(s) \, ds < \infty. \label{e4.28}
\end{equation}
As in \cite{MM181}, we take
\begin{align*}
J(t) := q \int_{t_0}^{t} \int_{\Omega} |\nabla u(t)
- \nabla u(t - s)|^2 \, dx \, ds ,
\end{align*}
where \eqref{e4.28} allows us to choose a constant $0 <q<1$ so that
\begin{equation}
J(t) < 1, \quad \forall  t \geq t_0. \label{e4.29}
\end{equation}
Then, we take
\[
\lambda(t) := - \int_{t_0}^{t} g'(s) \int_{\Omega} |\nabla u(t)
- \nabla u(t - s)|^2 \, dx \, ds
\]
and observe that $ \lambda(t) \leq - c E'(t)$.
Since $G$ is strictly convex on $(0, r]$ and $G(0) =0$, it follows that
\[
G(\theta z) \leq \theta G(z),\quad \text{for }  0 \leq \theta \leq 1  \text{ and }
z \in (0, r].
\]
By employing the above estimate, hypothesis (A1), \eqref{e4.29} and Jensen's inequality,
we obtain
\begin{align*}
\lambda(t)
&= \frac{1}{J(t)} \int_{t_0}^{t} J(t)(-g'(s)) \int_{\Omega} \left(|\nabla u(t)
 - \nabla u(t - s)|^2\right) \, dx \, ds \\
&\geq \frac{1}{J(t)} \int_{t_0}^{t} J(t) \xi(s) G(g(s))
 \int_{\Omega} \left(|\nabla u(t) - \nabla u(t - s)|^2 \right) \, dx \, ds \\
&\geq \frac{\xi(t)}{J(t)} \int_{t_0}^{t} G \left(J(t)g(s)\right)
 \int_{\Omega} \left(|\nabla u(t) - \nabla u(t - s)|^2 \right) \, dx \, ds \\
&\geq \frac{\xi(t)}{q} G\Big(\frac{1}{J(t)} \int_{t_0}^{t} J(t)g(s) q
\int_{\Omega} \left(|\nabla u(t) - \nabla u(t - s)|^2 \right) \, dx \, ds \Big) \\
&= \frac{\xi(t)}{q} G\Big(q\int_{t_0}^{t} g(s) \int_{\Omega}
\left(|\nabla u(t) - \nabla u(t - s)|^2\right) \, dx \, ds \Big) \\
&= \frac{\xi(t)}{q} \overline{G} \Big(q \int_{t_0}^{t} g(s)
\int_{\Omega} \left(|\nabla u(t) - \nabla u(t - s)|^2 \right) \, dx \, ds \Big),
\end{align*}
where $\overline{G}$ is an extension of $G$ such that $\overline{G}$ is strictly
increasing and strictly convex $C^2$ function on $(0, \infty)$.
This indicates that
\begin{align*}
 \int_{t_0}^t g(s) \int_{\Omega} \left(|\nabla u(t) - \nabla u(t - s)|^2 \right) \, dx
  \, ds
  \leq \frac{1}{q} \overline{G}^{-1} \Big( \frac{q \lambda(t)}{\xi(t) } \Big),
\end{align*}
and \eqref{e4.27} becomes
\begin{equation}
F'(t) \leq -k E(t) + \frac{c}{q}  \overline{G}^{-1}
\Big( \frac{q \lambda(t)}{\xi(t)} \Big), \quad \forall t \geq t_0.\label{e4.30}
\end{equation}
For $\varepsilon_0 < r$, using \eqref{e4.4}, and the fact that $E' \leq 0$,
$\overline{G}' > 0$, $\overline{G}'' > 0$, we find that the functional $F_1$,
defined by
\[
F_1(t) := \overline{G}' \Big( \frac{\varepsilon_0 E(t)}{E(0)} \Big) F(t) ,
\]
satisfies
\begin{equation}
\begin{aligned}
F_1'(t)
& = \frac{\varepsilon_0 E'(t)}{E(0)} \overline{G}''
\Big( \frac{\varepsilon_0 E(t)}{E(0)} \Big) F(t)
 + \overline{G}' \Big( \frac{\varepsilon_0 E(t)}{E(0)} \Big) F'(t) \\
& \leq -k E(t)\overline{G}' \Big( \frac{\varepsilon_0 E(t)}{E(0)} \Big)
 + \frac{c}{q} \overline{G}' \Big( \frac{\varepsilon_0 E(t)}{E(0)} \Big)
 \overline{G}^{-1} \Big( \frac{q\lambda(t)}{\xi(t)} \Big).
\end{aligned} \label{e4.31}
\end{equation}
Consider $\overline{G}^{*}$ the convex conjugate of $\overline{G}$ in the sense of
Young \cite[pp. 61-64]{AkilM}, given by
\begin{equation}
\overline{G}^{*}= s (\overline{G}')^{-1}(s) - \overline{G}[(\overline{G}')^{-1}(s)], \label{e4.32}
\end{equation}
and which satisfies the  generalized Young inequality
\begin{equation}
AB \leq \overline{G}^{*}(A) + \overline{G}(B). \label{e4.33}
\end{equation}
Then, with $A= \overline{G}'\big( \frac{\varepsilon_0 E(t)}{E(0)} \big)$ and
$B = \overline{G}^{-1} \big( \frac{q\lambda(t)}{\xi(t)} \big)$, and using
\eqref{e3.1} and \eqref{e4.5}--\eqref{e4.7}, we attain
\begin{align*}
F'_1(t)
&\leq -k E(t) \overline{G}'\Big(\frac{\varepsilon_0 E(t)}{E(0)}\Big)
+ \frac{c}{q}\overline{G}^{*}\Big( \overline{G}'
 \Big(\frac{\varepsilon_0 E(t)}{E(0)}\Big) \Big)
+ \frac{c}{q} \frac{q\lambda(t)}{\xi(t)} \\
&\leq -k E (t)\overline{G}'\Big(\frac{\varepsilon_0 E(t)}{E(0)}\Big)
+ \frac{c}{q} \varepsilon_0 \frac{E(t)}{E(0)} \overline{G}'
 \Big(\frac{\varepsilon_0 E(t)}{E(0)}\Big)
 + \frac{c \lambda(t)}{\xi(t)}
\\
&\leq -k E (t)\overline{G}'\Big(\frac{\varepsilon_0 E(t)}{E(0)}\Big)
+ \widehat{c}\varepsilon_0 \frac{E(t)}{E(0)}
 \overline{G}'\Big(\frac{\varepsilon_0 E(t)}{E(0)}\Big)
 + \frac{c \lambda(t)}{\xi(t)}\\
&= -\left(k E(0)- \widehat{c}\varepsilon_0 \right)\frac{E(t)}{E(0)}
 \overline{G}'\Big(\frac{\varepsilon_0 E(t)}{E(0)}\Big) + \frac{c \lambda(t)}{\xi(t)}.
\end{align*}
After fixing $\varepsilon_0$ much smaller (if needed), we arrive at
\begin{align*}
F'_1(t) \leq - k_1\frac{E(t)}{E(0)}
\overline{G}'\Big(\frac{\varepsilon_0 E(t)}{E(0)}\Big) + \frac{c \lambda(t)}{\xi(t)},
\end{align*}
where $k_1>0$.

Then, we multiply by $\xi(t)$ and take $\varepsilon_0$ so small that,
$\varepsilon_0 \frac{E(t)}{E(0)} < r$,
$\overline{G}'\Big(\varepsilon_0 \frac{E(t)}{E(0)}\Big)
= G'\Big(\varepsilon_0 \frac{E(t)}{E(0)}\Big)$ to arrive at
\begin{align*}
\xi(t)F'_1(t)
& \leq -k_1 \xi(t) \frac{E(t)}{E(0)} G'\Big(\varepsilon_0 \frac{E(t)}{E(0)}\Big)
+ cq\lambda(t) \\
& \leq -k_1\xi(t)\frac{E(t)}{E(0)}G'\Big(\varepsilon_0 \frac{E(t)}{E(0)}\Big)
- c E'(t).
\end{align*}
Let $F_2 = \xi F_1 + c E$ and recall that $\xi$ is non-increasing, to obtain
\[
F'_2(t)\leq \xi(t) F_1'(t) + c E'(t)
 \leq -k_1 \xi(t) \frac{E(t)}{E(0)} G'\Big(\varepsilon_0 \frac{E(t)}{E(0)}\Big),
\]
or
\begin{equation}
 k_1 \xi(t) \frac{E(t)}{E(0)} G'\Big(\varepsilon_0 \frac{E(t)}{E(0)}\Big)
 \leq -F'_2(t). \label{e4.34}
\end{equation}
Using the strictly increasing property of $G'$, non-increasing properties
of $E$ and $\xi$, we obtain for $t \geq t_0$,
 \begin{align*}
 k_1 \frac{\varepsilon_0 E(t)}{E(0)} G'\Big(\varepsilon_0 \frac{E(t)}{E(0)}\Big)
 \int_{t_0}^t \xi(s)ds \leq \varepsilon_0 F_2(t_0).
\end{align*}
Next, we set $ G_2(\tau ) = \tau G'(\tau )$ which is strictly increasing,
then we obtain, for some fixed positive constants $C $ and $\gamma$
\begin{align*}
E(t)\leq C G_2^{-1}\Big( \frac{\gamma}{\int_{t_0}^t \xi(s)ds}\Big), \quad
\forall t \geq t_0.
\end{align*}
This complete the proof.
\end{proof}
% Source-native exact authored range 1649--1793
\begin{thebibliography}{10}
\bibitem{Assila1} M. Aassila;
\emph{A note on the boundary stabilization of a compactly coupled system of wave
equations}.
 Applied Mathematics Letters 12.3 (1999): 19-24.

\bibitem{Assila2}  Mohammed Aassila;
\emph{Strong asymptotic stability of a compactly coupled system of wave equations}.
Applied Mathematics Letters 14.3 (2001): 285-290.

\bibitem{AkilM}  Mohammad, Akil, et al.;
\emph{Stability results of coupled wave models with locally memory in a past history
framework via nonsmooth coefficients on the interface}.
 Mathematical Methods in the Applied Sciences 44.8 (2021): 6950-6981.

\bibitem{Beniani} Abderrahmane Beniani,  Noureddine Taouaf,  Abbes Benaissa;
\emph{Well-posedness and exponential stability for coupled Lam\'e system with
viscoelastic term and strong damping}. Computers Mathematics with Applications 75.12 (2018): 4397-4404.

\bibitem{Berrimi2004}  Said Berrimi, Salim A. Messaoudi;
\emph{Exponential decay of solutions to a viscoelastic equation with nonlinear 
localized damping}. Electronic Journal of Differential Equations,
 2004 (2004), No. 88, 1=10.
 
\bibitem{Berrimi2006} Said Berrimi, Salim A. Messaoudi;
\emph{Existence and decay of solutions of a viscoelastic equation with a 
nonlinear source}. Nonlinear Analysis: Theory, Methods Applications 64.10 
(2006): 2314-2331.

\bibitem{CavalcantiDS} M. M. Cavalcanti, VN Domingos Cavalcanti, J. Ferreira;
\emph{Existence and uniform decay for a non-linear viscoelastic equation with 
strong damping}. Mathematical methods in the applied sciences 24.14 (2001): 1043-1053.

\bibitem{CavalcantiFerreira} M. M. Cavalcanti, VN Domingos Cavalcanti, J. Ferreira;
\emph{Exponential decay for the solution of semilinear viscoelastic wave equation 
with localized damping}. Electronic J. Diff. Eqns., 2002 (2002) No. 44, 1-14.

\bibitem{Cavalcanti2003} M. M. Cavalcanti, Higidio Portillo Oquendo;
\emph{Frictional versus viscoelastic damping in a semilinear wave equation}.
 SIAM journal on control and optimization 42.4 (2003): 1310-1324.
 
\bibitem{Dafermos1} Constantine M. Dafermos;
\emph{An abstract Volterra equation with applications to linear viscoelasticity}.
Journal of Differential Equations 7.3 (1970): 554-569.

\bibitem{Dafermos2} Constantine M. Dafermos;
\emph{Asymptotic stability in viscoelasticity}.
Archive for rational mechanics and analysis 37 (1970): 297-308.

\bibitem{FengBaoweiHaiyan} Baowei Feng, Haiyan Li;
\emph{Decay rates for a coupled viscoelastic Lam\'e system with strong damping}.
Mathematical Modelling and Analysis 25.2 (2020): 226-240.

\bibitem{GASM} Aissa Guesmia, Salim A. Messaoudi, Mostafa Zahri;
\emph{General decay of solutions of a weakly coupled abstract evolution equations 
with one finite memory control}.
 Journal of Applied Analysis, Computation 14.6 (2024): 3539-3557.
 
\bibitem{Hao} Jianghao Hao,  Fei Wang;
\emph{General decay rate for weak viscoelastic wave equation with Balakrishnan-Taylor 
damping and time-varying delay}.
Computers Mathematics with Applications 78.8 (2019): 2632-2640.

\bibitem{JKJT} Kun-Peng Jin, Jin Liang,  Ti-Jun Xiao;
\emph{Coupled second order evolution equations with fading memory: Optimal energy
 decay rate}. Journal of Differential Equations 257.5 (2014): 1501-1528.
 
\bibitem{KVBR} Vilmos Komornik, Bopeng Rao;
\emph{Boundary stabilization of compactly coupled wave equations}.
 Asymptotic Analysis 14.4 (1997): 339-359.
 
\bibitem{Li2011} Fushan Li, Zengqin Zhao, Yanfu Chen;
\emph{Global existence uniqueness and decay estimates for nonlinear viscoelastic 
wave equation with boundary dissipation}.
 Nonlinear Analysis: Real World Applications 12.3 (2011): 1759-1773.
 
\bibitem{Lions} J. L. Lions;
\emph{Quelques methodes de resolution des problemes aux limites non lineaires}.
 Dunford/Gauthier-Villars (1969).
 
\bibitem{LRF} Renato F.C. Lobato, et al.;
\emph{Optimal polynomial decay to coupled wave equations and its numerical properties}
 Journal of Applied Mathematics 2014.1 (2014): 897080.
 
\bibitem{SMessaoudi1} Salim A. Messaoudi;
\emph{General decay of solutions of a viscoelastic equation}.
Journal of Mathematical Analysis and Applications 341.2 (2008): 1457-1467.

\bibitem{SMessaoudi2} Salim Messaoudi;
\emph{General decay of the solution energy in a viscoelastic equation with a nonlinear 
source}. Nonlinear Analysis: Theory, Methods Applications 69.8 (2008): 2589-2598.

\bibitem{MessaoudiMostafa} Salim A. Messaoudi, Mostafa Zahri;
\emph{Analytical and computational results for the decay of solutions of a damped 
wave equation with variable-exponent nonlinearities}. (2022): 851-866.

\bibitem{MessaoudiMustafa} Salim A. Messaoudi,  Muhammad Islam Mustafa;
\emph{On the control of solutions of viscoelastic equations with boundary feedback}.
 Nonlinear Analysis: Real World Applications 10.5 (2009): 3132-3140.
 
\bibitem{MessaoudiWaled} Salim A. Messaoudi, Waled Al-Khulaifi;
\emph{General and optimal decay for a quasilinear viscoelastic equation}.
Applied Mathematics Letters 66 (2017): 16-22.

\bibitem{MTatar1} Salim A. Messaoudi,  N. E. Tatar;
\emph{Global existence and asymptotic behavior for nonlinear viscoelastic problem}.
 Mathematical Sciences Research Journal 7 (2003): 136-149.
 
\bibitem{MTatar2} Salim A. Messaoudi,  Nasser-eddine Tatar;
\emph{Global existence and uniform stability of solutions for a quasilinear 
viscoelastic problem}. Mathematical Methods in the Applied Sciences 30 (6) (2007).

\bibitem{MM181} Muhammad I. Mustafa;
\emph{Optimal decay rates for the viscoelastic wave equation}.
 Mathematical Methods in the Applied Sciences 41.1 (2018): 192-204.
 
\bibitem{MM} Muhammad I. Mustafa, Salim A. Messaoudi;
\emph{General stability result for viscoelastic wave equations}.
 Journal of Mathematical Physics 53.5 (2012).
 
\bibitem{MustafaMuhammadR} Muhammad I. Mustafa;
\emph{General decay result for nonlinear viscoelastic equations}.
 Journal of Mathematical Analysis and Applications 457.1 (2018): 134-152.
 
 \bibitem{Munoz} Jaime E. Munoz Rivera, Doherty Andrade;
\emph{Exponential decay of non-linear wave equation with a viscoelastic boundary 
condition}. Mathematical Methods in the Applied Sciences 23.1 (2000): 41-61.


\bibitem{ORHC} Rafael Lima Oliveira, Higidio Portillo Oquendo, Celene Buriol;
\emph{Exact decay rates for weakly coupled systems with indirect damping by 
fractional memory effects}. Mathematische Nachrichten 296.10 (2023): 4610-4633.

\bibitem{PJS} Jong Yeoul Park, Sun Hye Park;
\emph{General decay for quasilinear viscoelastic equations with nonlinear weak damping}.
 Journal of Mathematical Physics 50.8 (2009).
 
\bibitem{SaidH2011}  Belkacem Said-Houari, Salim A. Messaoudi, Aissa Guesmia;
\emph{General decay of solutions of a nonlinear system of viscoelastic wave equations}.
Nonlinear Differential Equations and Applications NoDEA 18.6 (2011): 659-684.

\bibitem{WSH} Shun-Tang Wu;
\emph{General decay for a wave equation of Kirchhoff type with a boundary control 
of memory type}. Boundary Value Problems (2011): 1-15.

\end{thebibliography}
\end{document}
